\documentclass[11pt]{article}
\usepackage[margin=1in]{geometry}
\usepackage{amsmath,amssymb,amsthm,booktabs,longtable,array,calc}
\usepackage[hidelinks]{hyperref}
\usepackage[T1]{fontenc}
\usepackage[utf8]{inputenc}
\newtheorem{theorem}{Theorem}[section]
\newtheorem{proposition}[theorem]{Proposition}
\newtheorem{lemma}[theorem]{Lemma}
\newtheorem{corollary}[theorem]{Corollary}
\theoremstyle{definition}
\newtheorem{definition}[theorem]{Definition}
\theoremstyle{remark}
\newtheorem*{remark}{Remark}
\providecommand{\tightlist}{\setlength{\itemsep}{0pt}\setlength{\parskip}{0pt}}
\begin{document}
\appendix
\section{\texorpdfstring{Reconstruction and
block-compatibility closure for
\(MVMMMM\)}{Reconstruction and block-compatibility closure for MVMMMM}}\label{appendix-e.-reconstruction-and-block-compatibility-closure-for-mvmmmm}

This appendix keeps the two logical ingredients of Theorem 7.2 separate:
fixed-row reconstruction/full validity, and the independent block
criterion.

\subsection{E.1. Fixed-row semantics}\label{e.1.-fixed-row-semantics}

For \(MVMMMM\), every horizontal relation is \[
A_c\prec B_c.
\] The upper row is \(R^+(k,a)\) and the lower row is \(R^-(k,b)\).
Hence the upper column sequence is the H-opener stream, while the lower
column sequence is the H-closer request stream.

\subsection{E.2. Fixed-row reconstruction: H/H noncrossing is
LIFO}\label{e.2.-fixed-row-reconstruction-hh-noncrossing-is-lifo}

With each opener before its matching closer, two horizontal intervals
cross exactly when an earlier-opened interval closes while a
later-opened distinct interval is still open. Pairwise \(H/H\) Butterfly
noncrossing is therefore equivalent to the usual LIFO stack rule.

Fix the two row chains and let \(c\) be the next requested lower label.
If the matching upper opener has not yet been emitted, no lower label
may bypass \(c\), so the upper chain is forced until that opener
appears. If \(c\) is the stack top, its closer is forced immediately;
pushing another opener would block the fixed next closer. If \(c\) is
already buried in the stack, no completion exists. Induction over the
lower events proves that an H-compatible merge is unique when it exists,
and that the forced reconstruction succeeds exactly when H compatibility
holds.

\subsection{E.3. Automatic remainder after horizontal
reconstruction}\label{e.3.-automatic-remainder-after-horizontal-reconstruction}

Proposition 6.1 supplies all upper and lower vertical directed
constraints and all same-row same-target vertical noncrossing. The H
stack supplies every horizontal precedence and every \(H/H\) noncrossing
condition. Since the horizontal target \(S\) is distinct from the
vertical targets \(E,W\), no mixed horizontal/vertical twin family
occurs. The only remaining class is \(U_i/L_j\) with equal seam parity.

For seam \(i\), orient the adjacent horizontal nesting as \[
A_{\alpha_i}\prec A_{\beta_i}\prec B_{\beta_i}\prec B_{\alpha_i},
\] where \((\alpha_i,\beta_i)=(i,i+1)\) for odd \(i\) and \((i+1,i)\)
for even \(i\). Thus \(U_i\) is the opener-side interval and \(L_i\) the
closer-side interval of the same nested H pair.

For distinct equal-parity seams \(i,j\), the two adjacent column pairs
are disjoint. Consider two nested H pairs \[
o_x<o_y<c_y<c_x,
\qquad
o_p<o_q<c_q<c_p,
\] with opener-side intervals \(U_e=[o_x,o_y]\), \(U_f=[o_p,o_q]\) and
closer-side intervals \(L_e=[c_y,c_x]\), \(L_f=[c_q,c_p]\). Suppose
\(U_e\) crosses \(L_f\).

If \[
c_q<o_x<c_p<o_y,
\] then \[
o_p<o_x<c_p<c_x,
\] so the horizontal intervals \([o_p,c_p]\) and \([o_x,c_x]\) cross,
contradicting horizontal laminarity. Hence the only possible alternating
order is \[
o_x<c_q<o_y<c_p.
\] If \(o_p<o_x\), horizontal laminarity of \([o_q,c_q]\) with
\([o_x,c_x]\) forces \(o_x<o_q\), giving \[
o_p<o_x<o_q<o_y,
\] so \(U_e\) crosses \(U_f\). If instead \(o_x<o_p\), laminarity of
\([o_p,c_p]\) with \([o_y,c_y]\) forces \(c_y<c_p\), while laminarity of
\([o_x,c_x]\) with \([o_p,c_p]\) forces \(c_p<c_x\). Together with
\(c_q<o_y<c_y\) this yields \[
c_q<c_y<c_p<c_x,
\] so \(L_e\) crosses \(L_f\). Therefore \[
U_i\text{ crosses }L_j
\Longrightarrow
(U_i\text{ crosses }U_j)\ \text{or}\ (L_i\text{ crosses }L_j).
\] Equal parity gives the same vertical target class, and Proposition
6.1 forbids both same-row alternatives. For \(i=j\), the upper interval
lies in the opener side of the nested adjacent H pair and the lower
interval in the closer side, so they are disjoint. Thus every
H-compatible merge is globally valid: \[
\operatorname{FullCompat}(k;a,b)
\iff
\operatorname{HCompat}(k;a,b).
\]

\subsection{E.4. Endpoint normal forms for the block
criterion}\label{e.4.-endpoint-normal-forms-for-the-block-criterion}

Let \(W_p\) denote the \(+\) row-normal-form column word for pivot
\(p\). Because the lower row is the reverse normal form, its closer
request stream is \[
\overleftarrow{W}_p:=\operatorname{rev}(W_p).
\] The singleton-block word is \[
W_1=(1,3,5,\ldots,2k-1,2k,2k-2,\ldots,2).
\] For \(1\le r\le k-1\), name the two pivots in block \(r\) by \[
p_r^{(0)}=2r,
\qquad
p_r^{(1)}=2r+1.
\] Set \(c_r=2r+1\), \(d_r=2r+2\) and \[
L_r=(2r-1,2r-3,\ldots,1,2,4,\ldots,2r),
\] \[
R_r=(2r+3,2r+5,\ldots,2k-1,2k,2k-2,\ldots,2r+4).
\] Then \[
W_{p_r^{(0)}}=c_rR_rd_rL_r,
\qquad
W_{p_r^{(1)}}=c_rL_rR_rd_r,
\] and \[
\overleftarrow{W}_{p_r^{(0)}}=
\operatorname{rev}(L_r)d_r\operatorname{rev}(R_r)c_r,
\] \[
\overleftarrow{W}_{p_r^{(1)}}=
d_r\operatorname{rev}(R_r)\operatorname{rev}(L_r)c_r.
\] At \(r=k-1\), \(R_r\) is empty and the formulas remain literal. The
block coordinate records the common junction index: \[
B_R(1)=0,
\qquad
B_R(p_r^{(0)})=B_R(p_r^{(1)})=r.
\]

\subsection{E.5. Necessity for the block criterion: unequal
blocks}\label{e.5.-necessity-for-the-block-criterion-unequal-blocks}

For three distinct horizontal labels \(x,y,z\), suppose the opener order
contains \[
y\prec_W z\prec_W x
\] while the closer stream requires \[
x\prec_{\leftarrow W} y\prec_{\leftarrow W} z.
\] No LIFO schedule can realize both. When \(x\) closes, both \(y\) and
\(z\) have already opened; after \(x\) is removed, \(z\) lies above
\(y\), yet the prescribed closer stream asks for \(y\) before \(z\).
Additional pushes cannot repair the order and \(z\) is not permitted to
close first.

Every unequal ordered block pair contains such a witness. Since
\(\overleftarrow{W}_b=\operatorname{rev}(W_b)\), it is enough to verify
\[
y\prec_{W_a}z\prec_{W_a}x,
\qquad
z\prec_{W_b}y\prec_{W_b}x.
\] The following table covers all ordered unequal-block cases directly.

\begin{longtable}[]{@{}
  >{\raggedright\arraybackslash}p{(\columnwidth - 8\tabcolsep) * \real{0.2000}}
  >{\raggedright\arraybackslash}p{(\columnwidth - 8\tabcolsep) * \real{0.2000}}
  >{\raggedright\arraybackslash}p{(\columnwidth - 8\tabcolsep) * \real{0.2000}}
  >{\raggedright\arraybackslash}p{(\columnwidth - 8\tabcolsep) * \real{0.2000}}
  >{\raggedright\arraybackslash}p{(\columnwidth - 8\tabcolsep) * \real{0.2000}}@{}}
\toprule\noalign{}
\begin{minipage}[b]{\linewidth}\raggedright
ordered case
\end{minipage} & \begin{minipage}[b]{\linewidth}\raggedright
range
\end{minipage} & \begin{minipage}[b]{\linewidth}\raggedright
\((x,y,z)\)
\end{minipage} & \begin{minipage}[b]{\linewidth}\raggedright
order in \(W_a\)
\end{minipage} & \begin{minipage}[b]{\linewidth}\raggedright
order in \(W_b\)
\end{minipage} \\
\midrule\noalign{}
\endhead
\bottomrule\noalign{}
\endlastfoot
\(1/p_s^{(0)}\) & \(1\le s\le k-1\) & \((2,1,2s+1)\) & \(1<2s+1<2\) &
\(2s+1<1<2\) \\
\(1/p_s^{(1)}\) & \(1\le s\le k-1\) & \((2,1,2s+1)\) & \(1<2s+1<2\) &
\(2s+1<1<2\) \\
\(p_r^{(0)}/1\) & \(1\le r\le k-1\) & \((2,2r+1,1)\) & \(2r+1<1<2\) &
\(1<2r+1<2\) \\
\(p_r^{(1)}/1\) & \(1\le r\le k-1\) & \((2,2r+1,1)\) & \(2r+1<1<2\) &
\(1<2r+1<2\) \\
\(p_r^{(0)}/p_s^{(0)}\) or \(p_r^{(0)}/p_s^{(1)}\) & \(r<s\) &
\((1,2r+1,2r+3)\) & \(2r+1<2r+3<1\) & \(2r+3<2r+1<1\) \\
\(p_r^{(1)}/p_s^{(0)}\) or \(p_r^{(1)}/p_s^{(1)}\) & \(r<s\) &
\((2s,1,2r+3)\) & \(1<2r+3<2s\) & \(2r+3<1<2s\) \\
\(p_r^{(0)}/p_s^{(0)}\) or \(p_r^{(1)}/p_s^{(0)}\) & \(r>s\) &
\((1,2r+1,2r-1)\) & \(2r+1<2r-1<1\) & \(2r-1<2r+1<1\) \\
\(p_r^{(0)}/p_s^{(1)}\) or \(p_r^{(1)}/p_s^{(1)}\) & \(r>s\) &
\((2r,2r+1,1)\) & \(2r+1<1<2r\) & \(1<2r+1<2r\) \\
\end{longtable}

All witnesses lie in range. If \(r<s\), then \(r\le k-2\) and
\(2r+3\le2k-1\); if \(r>s\), then \(r\ge2\) and \(2r-1\ge3\);
singleton/non-singleton cases require only \(k\ge2\). Hence every
unequal-block pair fails H compatibility.

\subsection{E.6. Sufficiency for the block criterion: same-block
schedules}\label{e.6.-sufficiency-for-the-block-criterion-same-block-schedules}

For the singleton pair \((1,1)\), push all openers in \(W_1\) and pop
them in reverse order.

Within a nonzero block \(r\), the diagonal pairs \[
(p_r^{(0)},p_r^{(0)}),
\qquad
(p_r^{(1)},p_r^{(1)})
\] again use push-all/pop-all. For \((p_r^{(0)},p_r^{(1)})\), use \[
\text{push }c_r,R_r,d_r;\quad
\text{pop }d_r,\operatorname{rev}(R_r);\quad
\text{push }L_r;\quad
\text{pop }\operatorname{rev}(L_r);\quad
\text{pop }c_r.
\] For \((p_r^{(1)},p_r^{(0)})\), use \[
\text{push }c_r,L_r;\quad
\text{pop }\operatorname{rev}(L_r);\quad
\text{push }R_r,d_r;\quad
\text{pop }d_r,\operatorname{rev}(R_r);\quad
\text{pop }c_r.
\] Every requested closer is the stack top, every label is used once,
and the stack ends empty. Thus equal block coordinate implies H
compatibility.

Combining E.5 and E.6 gives \[
\operatorname{HCompat}(k;a,b)
\iff
B_R(a)=B_R(b).
\] Together with E.3 this proves Theorem 7.2. No finite compatibility
graph or count formula is used as a premise.

\end{document}
